% from main.tex
%%%%%%%%%%  scan idx 109  |  folio 96  %%%%%%%%%%
\textbf{Proposition 2.2.} \emph{Soient \(F\in\Ob\Dpx(\mathcal C_S)\) et \(n\in\mathbb Z\). Alors} :

a) \emph{Si \(F\) est strictement \(n\)-pseudo-cohérent (relativement à \(\mathcal C_0\)) et si \(u:E\to F\) est un quasi-isomorphisme, l'ensemble des \(X\) au-dessus de \(S\) tels qu'il existe un quasi-isomorphisme \(v:F'\longrightarrow E_X\) avec \(F'\) strictement \(n\)-pseudo-cohérent est un raffinement de \(S\).}

b) \emph{Les conditions suivantes sont équivalentes} :
\begin{enumerate}[label=(\roman*)]
\item \emph{l'ensemble des \(X\) au-dessus de \(S\) tels qu'il existe un quasi-isomorphisme \(u:E\longrightarrow F_X\) de \(\Cpx(\mathcal C_X)\), avec \(E\) strictement \(n\)-pseudo-cohérent, est un raffinement de \(S\)}~;
\item[\((\mathrm{i}')\)] \emph{même condition que (i), mais \(u\) étant un isomorphisme de \(\Dpx(\mathcal C_X)\)}~;
\item \emph{l'ensemble des \(X\) au-dessus de \(S\) tels qu'il existe un \(n\)-quasi-isomorphisme (1.4.3) \(u:E\longrightarrow F_X\) de \(\Cpx(\mathcal C_X)\), avec \(E\) strictement parfait, est un raffinement de \(S\)}~;
\item[\((\mathrm{ii}')\)] \emph{même condition que (ii), mais \(u\) étant un \(n\)-isomorphisme de \(\Dpx(\mathcal C_X)\)}.
\end{enumerate}

\emph{Démonstration.} a) Il existe, pour un entier \(i\) assez grand, un \(i\)-quasi-isomorphisme \(v_i:F'_i\longrightarrow E\), avec \(F'_i\) strictement \(n\)-pseudo-cohérent (prendre \(F'_i=0\)). Montrons que, si \(i>n\), il existe, quitte à raffiner \(S\), un \((i-1)\)-quasi-isomorphisme \(v_{i-1}:F'_{i-1}\longrightarrow E\), avec \(F'_{i-1}\) strictement \(n\)-pseudo-cohérent. D'après (1.4.4), il suffit de montrer que \(\Ho^{i-1}(M')\), où \(M'\) est le cône de \(v_i\), est de \(\mathcal C_0\)-type fini. Or, si \(M''\) désigne le cône de \(uv_i\), l'hypothèse que \(u\) est un quasi-isomorphisme et (1.4.2) impliquent que l'on a
\[
  \Ho^j(M')\qis \Ho^j(M'')\qquad\text{pour tout }j
\]
et \(\Ho^j(M')=0\) pour \(j\geq i\).
